% =====================================================================
%  Publisher Submission Sheet — Book Cover & Manuscript Metadata
%  For: "The Blind Leading the Blind" (monograph)
%  Author: Kundai Farai Sachikonye
%
%  Each field is annotated with its character count against the
%  publisher's limit so the values can be verified at a glance.
% =====================================================================
\documentclass[11pt,a4paper]{article}

\usepackage[utf8]{inputenc}
\usepackage[T1]{fontenc}
\usepackage{lmodern}
\usepackage[a4paper,margin=2.4cm]{geometry}
\usepackage{microtype}
\usepackage{xcolor}
\usepackage{booktabs}
\usepackage{enumitem}
\usepackage[colorlinks=true,linkcolor=black,urlcolor=blue!50!black]{hyperref}

\definecolor{fieldbg}{RGB}{245,246,248}
\definecolor{rulecol}{RGB}{30,64,120}
\definecolor{labelcol}{RGB}{30,64,120}
\definecolor{limitcol}{RGB}{120,120,120}

% ---- a labelled field block -----------------------------------------
\newcommand{\field}[3]{% {label}{limit note}{content}
  \par\medskip
  {\color{labelcol}\bfseries\large #1}%
  \hfill{\color{limitcol}\small #2}\par
  \vspace{2pt}%
  \noindent\colorbox{fieldbg}{%
    \begin{minipage}{\dimexpr\linewidth-2\fboxsep\relax}
    \vspace{2pt}#3\vspace{2pt}
    \end{minipage}}%
  \par
}

\newcommand{\charcount}[2]{{\color{limitcol}\small\itshape [#1 characters / #2 limit]}}
\newcommand{\wordcount}[1]{{\color{limitcol}\small\itshape [#1 words]}}

% Section heading with an underrule (no titlesec dependency):
\newcommand{\Section}[1]{%
  \par\addvspace{16pt}%
  {\color{rulecol}\Large\bfseries #1}\par
  \vspace{2pt}{\color{rulecol}\rule{\linewidth}{1pt}}\par\vspace{4pt}}

\setlength{\parindent}{0pt}
\setlength{\parskip}{2pt}

\begin{document}

% =====================================================================
\begin{center}
  {\color{rulecol}\Huge\bfseries Publisher Submission Sheet}\\[6pt]
  {\large Book Cover \& Manuscript Metadata}\\[10pt]
  {\color{limitcol}\rule{\linewidth}{0.8pt}}
\end{center}

\vspace{4pt}

% =====================================================================
\Section{Book Cover}

\field{Main Author Name}{}{%
  Kundai Farai Sachikonye
}

\field{Co-Authors' Names}{max.\ 2}{%
  \textit{None --- single-authored monograph.}
}

% Title: "The Blind Leading the Blind" = 27 chars
\field{Book Title}{100-character limit}{%
  The Blind Leading the Blind
  \par\vspace{3pt}\charcount{27}{100}
}

% Subtitle chosen to fit 100 chars while carrying the scientific content.
% "Categorical Resolution of Biological Dynamics: Cellular State, Disease, and Therapy"
%  count = 82 characters
\field{Subtitle}{100-character limit}{%
  Categorical Resolution of Biological Dynamics: Cellular State, Disease, and Therapy
  \par\vspace{3pt}\charcount{83}{100}
}

% Spine: "The Blind Leading the Blind  ·  Sachikonye" = 42 chars
\field{Spine}{60-character limit}{%
  The Blind Leading the Blind \quad Sachikonye
  \par\vspace{3pt}\charcount{40}{60}
}

% =====================================================================
% BLURB — target 250+ words, <=2200 chars. The version below is ~1980
% characters including spaces, ~330 words.
% =====================================================================
\field{Blurb}{2200-character limit; min.\ 250 words}{%
In 1568 Pieter Bruegel painted six blind men falling into a ditch, each with a
different, clinically diagnosable eye disease. This monograph takes that image
as a theorem about biology: every finite observer---every instrument, every
receptor, every cell---is blind in its own specific way, and no ensemble of
blind observers corrects itself by aggregation. Yet a healthy body still
navigates. How?

From just two axioms---that physical systems occupy a bounded phase space, and
that any finite observer partitions it into distinguishable cells---the book
derives a complete, self-contained mathematical framework for cellular state,
disease, and therapy. It proves that every bounded receiver has a strictly
positive epistemic ``floor'': the truth about a state is never a point but a
region of positive measure. From this follow the No Template Theorem (no cell
can store a model of itself, so self-diagnosis is structurally impossible), the
Group Blindness Theorem (collective observation cannot reach zero error), and
the Local Invisibility Theorem (disease is invisible to any observer confined
to a single reaction).

The resolution is geometric. Health is not knowledge of health but
\emph{self-consistency}: when the reaction rates around every metabolic loop
close on themselves---zero holonomy---the circuit is well, and this can be
checked from the inside with no external template. Therapy becomes the
restoration of loop closure, computed as a sparse, polynomial-time
optimisation. Purpose itself appears as a dynamical fixed point: the
destination is enacted, not seen.

Rigorous yet vivid, moving from Renaissance painting to category theory,
differential topology, and convex optimisation, this is a first-principles
account of why living systems work despite being built entirely from blind
parts.
  \par\vspace{4pt}\charcount{1799}{2200}\quad\wordcount{266}
}

% =====================================================================
% AUTHOR BIOGRAPHY — target 70+ words, <=500 chars. Below is ~470 chars, ~78 words.
% =====================================================================
\field{Author's Biography}{500-character limit; min.\ 70 words}{%
Kundai Farai Sachikonye is a researcher in bioinformatics at the Technical
University of Munich and the AIMe Registry for Artificial Intelligence. He
develops first-principles mathematical frameworks that unify observation,
computation, and physical dynamics, focused on the categorical geometry of
cellular state, disease, and therapy. His research asks how living systems stay
well while built entirely from bounded, error-prone parts. This present book is
the synthesis of that programme of work.
  \par\vspace{4pt}\charcount{499}{500}\quad\wordcount{70}
}

% =====================================================================
\Section{Manuscript}

\field{Book Language}{}{%
  English (British)
}

\field{Book Category}{}{%
  Science \guillemotright{} Life Sciences \guillemotright{} Biophysics \& Mathematical Biology
  \par\vspace{2pt}
  \textit{Secondary:} Mathematics \guillemotright{} Applied Mathematics (Dynamical Systems, Optimisation);\;
  Medicine \guillemotright{} Systems Biology \& Theoretical Medicine
}

\field{Keywords}{}{%
  \begin{minipage}{\linewidth}
  \begin{itemize}[leftmargin=1.4em,itemsep=1pt,topsep=2pt]
    \item categorical dynamics; partition geometry; bounded observers
    \item epistemic floor; cell-truth; resolution limit
    \item metabolic circuits; Wegscheider holonomy; loop closure
    \item disease as coherence deficit; No Template Theorem
    \item self-consistency; reference-free diagnosis
    \item sparse $\ell_1$ therapeutic optimisation; drug combination
    \item oscillator coherence; Kuramoto synchronisation
    \item mathematical biology; theoretical medicine; systems biology
  \end{itemize}
  \end{minipage}
}

\vfill
\begin{center}
{\color{limitcol}\rule{\linewidth}{0.8pt}}\\[3pt]
{\color{limitcol}\small Prepared for publisher submission \quad\textbullet\quad
Author: Kundai Farai Sachikonye \quad\textbullet\quad
\href{mailto:kundai.sachikonye@bitspark.com}{kundai.sachikonye@bitspark.com}}
\end{center}

\end{document}
